%
% 42-matrix.tex -- Matrixdarstellung algebraischer Elemente
%
% (c) 2026 Prof Dr Andreas Müller
%
\subsection{Matrixdarstellung algebraischer Elemente}
Für die Rechnung in einem Erweiterungskörper $\mathbb{Q}(\alpha)$
mit dem Minimalpolynom $m(X)\in \mathbb{Q}[X]$ ist bis jetzt kein
einfacher Algorithmus gegeben worden.
Die Matrixdarstellung eines algebraischen Elementes, die in diesem
Abschnitt studiert werden soll, liefert so einen Algorithmus.
Gleichzeitig wird damit aber auch ein Werkzeug bereitgestellt,
welches beim Beweis des Liouville-Prinzips für algebraische Erweiterungen
\index{Liouville-Prinzip}%
im Abschnitt~\ref{buch:intalg:liouville:subsection:allgemein}
als nützlich erweisen wird.

%
% Die Matrix eines algebraischen Elementes
%
\subsubsection{Die Matrix eines algebraischen Elements}
In diesem Abschnitt betrachten wir die Multiplikationsoperation
mit $\alpha$ im Erweiterungskörper $\mathbb{Q}(\alpha)$ für ein
algebraisches Element mit dem Minimalpolynom
\[
m(X)
=
m_0+m_1X+\dots + m_{n-1}X^{n-1}+X^n
\in
\mathbb{Q}[X].
\]
Der Erweiterungskörper ist die Menge
\[
\mathbb{Q}(\alpha)
=
\{
a_0+a_1\alpha+\ldots+a_{n-1}\alpha^{n-1}
\mid
a_k\in\mathbb{Q}\;
\forall k
\}.
\]
Er ist ein $n$-dimensionaler Vektorraum über $\mathbb{Q}$ mit der
Basis
\[
\mathscr{B}
=
\{
1,\alpha,\alpha^2,\ldots,\alpha^{n-1}
\}.
\]
Die Multiplikation mit $\alpha$ ist wegen
\[
\alpha (t u + s v)
=
t
\alpha
u
+
s
\alpha
v
,\qquad t,s\in\mathbb{Q},\;u,v\in\mathbb{Q}(\alpha)
\]
eine lineare Abbildung, man kann sie also auch in der Basis
$\mathscr{B}$ durch eine $n\times n$-Matrix $A$ beschreiben.
Aus der Abbildung der Basiselemente in $\mathscr{B}$ kann
man die Matrix
\begin{equation}
\left.
\begin{aligned}
1           &\mapsto \alpha   \\
\alpha      &\mapsto \alpha^2 \\
\alpha^2    &\mapsto \alpha^3 \\
\alpha^3    &\mapsto \alpha^4 \\[-3pt]
            &\phantom{n}\vdots\\
\alpha^{n-1}&\mapsto -a_0-a_1\alpha-\dots-a_{n-1}\alpha^{n-1}
\end{aligned}
\right\}
\quad\Rightarrow\quad
\bm{\alpha}
=
\left(
\raisebox{-1.39cm}{\begin{tikzpicture}[>=latex,thick]
\def\punkt#1#2{ ({(#1)*0.75},{-(#2)*0.5}) }
\clip \punkt{-0.5}{-0.5} rectangle \punkt{5.5}{5.5};
%\draw[color=darkred] \punkt{-0.5}{-0.5} rectangle \punkt{5.5}{5.5};
\def\element#1#2#3{
\node at \punkt{#1}{#2} {$#3\mathstrut$};
}
\def\horizontaldots#1#2{
\fill \punkt{#1-0.333}{#2} circle[radius=0.02];
\fill \punkt{#1}{#2} circle[radius=0.02];
\fill \punkt{#1+0.333}{#2} circle[radius=0.02];
}
\def\verticaldots#1#2{
\fill \punkt{#1}{#2-0.333} circle[radius=0.02];
\fill \punkt{#1}{#2} circle[radius=0.02];
\fill \punkt{#1}{#2+0.333} circle[radius=0.02];
}
\def\diagonaldots#1#2{
\fill \punkt{#1-0.333}{#2-0.333} circle[radius=0.02];
\fill \punkt{#1}{#2} circle[radius=0.02];
\fill \punkt{#1+0.333}{#2+0.333} circle[radius=0.02];
}
\draw[color=darkred!10,line width=11pt] \punkt{-1}{-1} -- \punkt{6}{6};
\draw[color=blue!10,line width=11pt] \punkt{-1}{0} -- \punkt{5}{6};
\element{0}{0}{0} \element{1}{0}{0} \element{2}{0}{0}
\horizontaldots{3}{0}
\element{4}{0}{0} \element{5}{0}{-m_0}
%
\element{0}{1}{1} \element{1}{1}{0} \element{2}{1}{0}
\horizontaldots{3}{1}
\element{4}{1}{0} \element{5}{1}{-m_1}
%
\element{0}{2}{0} \element{1}{2}{1} \element{2}{2}{0}
\horizontaldots{3}{2}
\element{4}{2}{0} \element{5}{2}{-m_2}
%
\element{0}{3}{0} \element{1}{3}{0} \element{2}{3}{1}
\diagonaldots{3}{3}
\element{5}{3}{-m_3}
%
\verticaldots{0}{4}
\verticaldots{1}{4}
\verticaldots{2}{4}
\diagonaldots{3}{4}
\diagonaldots{4}{4}
\verticaldots{5}{4}
%
\element{0}{5}{0} \element{1}{5}{0} \element{2}{5}{0}
\horizontaldots{3}{5}
\element{4}{5}{1} \element{5}{5}{-m_{n-1}}
\end{tikzpicture}}
\right)
\label{buch:komplex:fundamental:eqn:alphaMatrix}
\end{equation}
ablesen.

%
% Multiplikative Inverse in Q(\alpha)
%
\subsubsection{Multiplikative Inverse in $\mathbb{Q}(\alpha)$}
Ein Element $x$ von $\mathbb{Q}(\alpha)$ schreiben wir als Spaltenvektor mit
den Einträgen $x_0,x_1,\dots,x_{n-1}$.
Die Multiplikation mit $\alpha$ wird durch die Matrix $\bm{\alpha}$ von
\eqref{buch:komplex:fundamental:eqn:alphaMatrix} beschrieben.
Die Multiplikation des Elements $a\in \mathbb{Q}(\alpha)$ mit $x$ wird
daher durch die Matrixmultiplikation mit der Matrix
\[
A
=
a_0I+a_1\bm{\alpha} + a_2\bm{\alpha}^2 + \ldots + a_{n-1}\bm{\alpha}^{n-1}
\]
wiedergegeben.
Um die Inverse $x=a^{-1}$ zu finden, muss man also das lineare
Gleichungssystem
\index{lineares Gleichungssystem}%
$ Ax=e_1 $ lösen, wobei $e_1$ der erste Standardbasisvektor ist.
\index{Standardbasisvektor}%
Man kann also das multiplikative inverse Element von $a$ als
\index{multiplikative Inverse}%
\index{Inverse, multiplikativ}%
\[
x = A^{-1}e_1
\]
finden.

%
% Das charakteristische Polynom der Matrix A
%
\subsubsection{Das charakteristische Polynom der Matrix $\bm{\alpha}$}
Das charakteristische Polynom von $\bm{\alpha}$ kann am einfachsten
durch Entwicklung nach der letzten Spalte bestimmt werden.
\index{Entwicklung}%
Die Minordeterminante zum Element $-m_k$ mit $k<n-1$ ist
\index{Minordeterminante}%
\[
\det \bm{\alpha}_{k+1,n}
=
\left|
\;
\raisebox{-2.2cm}{\begin{tikzpicture}[>=latex,thick]
\def\punkt#1#2{ ({(#1)*0.75},{-(#2)*0.5}) }
\clip \punkt{-0.5}{-0.5} rectangle \punkt{8.5}{8.5};
\begin{scope}
	\clip \punkt{-0.5}{-0.5} rectangle \punkt{4.5}{4.5};
	\draw[color=darkred!10,line width=12pt] \punkt{-1}{-1} -- \punkt{4.5}{4.5};
	\draw[color=blue!10,line width=12pt] \punkt{-1}{0} -- \punkt{4.5}{5.5};
\end{scope}
\begin{scope}
	\clip \punkt{9.5}{9.5} rectangle \punkt{4.5}{4.5};
	\draw[color=blue!10,line width=12pt] \punkt{9}{9} -- \punkt{4.5}{4.5};
	\draw[color=darkred!10,line width=12pt] \punkt{9}{8} -- \punkt{4.5}{3.5};
\end{scope}
\draw[line width=0.4pt] \punkt{-0.5}{4.5} -- \punkt{8.5}{4.5};
\draw[line width=0.4pt] \punkt{4.5}{-0.5} -- \punkt{4.5}{8.5};

%\draw[color=darkred] \punkt{-0.5}{-0.5} rectangle \punkt{8.5}{8.5};
\def\element#1#2#3{
\node at \punkt{#1}{#2} {$#3\mathstrut$};
}
\def\horizontaldots#1#2{
\fill \punkt{#1-0.333}{#2} circle[radius=0.02];
\fill \punkt{#1}{#2} circle[radius=0.02];
\fill \punkt{#1+0.333}{#2} circle[radius=0.02];
}
\def\verticaldots#1#2{
\fill \punkt{#1}{#2-0.333} circle[radius=0.02];
\fill \punkt{#1}{#2} circle[radius=0.02];
\fill \punkt{#1}{#2+0.333} circle[radius=0.02];
}
\def\diagonaldots#1#2{
\fill \punkt{#1-0.333}{#2-0.333} circle[radius=0.02];
\fill \punkt{#1}{#2} circle[radius=0.02];
\fill \punkt{#1+0.333}{#2+0.333} circle[radius=0.02];
}
\element{0}{0}{-\lambda} \element{1}{0}{0}        \element{2}{0}{0}
\horizontaldots{3}{0}
\element{4}{0}{0}        \element{5}{0}{0}        \element{6}{0}{0}
\horizontaldots{7}{0}
\element{8}{0}{0}
%
\element{0}{1}{1}        \element{1}{1}{-\lambda} \element{2}{1}{0}
\horizontaldots{3}{1}
\element{4}{1}{0}        \element{5}{1}{0}        \element{6}{1}{0}
\horizontaldots{7}{1}
\element{8}{1}{0}
%
\element{0}{2}{0}        \element{1}{2}{1}        \element{2}{2}{-\lambda}
\horizontaldots{3}{2}
\element{4}{2}{0}        \element{5}{2}{0}        \element{6}{2}{0}
\horizontaldots{7}{2}
\element{8}{2}{0}
%
\verticaldots{0}{3}
\verticaldots{1}{3}
\verticaldots{2}{3}
\diagonaldots{3}{3}
\verticaldots{4}{3}
\verticaldots{5}{3}
\verticaldots{6}{3}
\diagonaldots{7}{3}
\verticaldots{8}{3}
%
\element{0}{4}{0}        \element{1}{4}{0}        \element{2}{4}{0}
\horizontaldots{3}{4}
\element{4}{4}{-\lambda} \element{5}{4}{0}        \element{6}{4}{0}
\horizontaldots{7}{4}
\element{8}{4}{0}
%
\element{0}{5}{0}        \element{1}{5}{0}        \element{2}{5}{0}
\horizontaldots{3}{5}
\element{4}{5}{0}        \element{5}{5}{1}        \element{6}{5}{-\lambda}
\horizontaldots{7}{5}
\element{8}{5}{0}
%
\element{0}{6}{0}        \element{1}{6}{0}        \element{2}{6}{0}
\horizontaldots{3}{6}
\element{4}{6}{0}        \element{5}{6}{0}        \element{6}{6}{1}
\horizontaldots{7}{6}
\element{8}{6}{0}
%
\verticaldots{0}{7}
\verticaldots{1}{7}
\verticaldots{2}{7}
\diagonaldots{3}{7}
\verticaldots{4}{7}
\verticaldots{5}{7}
\verticaldots{6}{7}
\diagonaldots{7}{7}
\verticaldots{8}{7}
%
\element{0}{8}{0} \element{1}{8}{0} \element{2}{8}{0}
\horizontaldots{3}{8}
\element{4}{8}{0} \element{5}{8}{0} \element{6}{8}{0}
\horizontaldots{7}{8}
\element{8}{8}{1}
\end{tikzpicture}}
%
\;
\right|
=
(-\lambda)^k.
\]
Damit tragen diese Unterdeterminanten insgesamt
\begin{align}
\sum_{k=0}^{n-2}
(-1)^{n+k}
\det \bm{\alpha}_{k+1,n}
&=
\sum_{k=0}^{n-2}
(-1)^{n+k}
(-m_k)
(-\lambda)^k
=
-
(-1)^n
\sum_{k=0}^{n-2}
m_k\lambda^k
\label{buch:komplex:fundamentalsatz:spur:eqn:unterdet}
\end{align}
bei.
Die Minordeterminante $\bm{\alpha}_{n,n}$ für das Element in der rechten
unteren Ecke lässt sich auf analoge Weise bestimmen, sie ist
$(-\lambda)^{n-1}$.
Sie leistet den Beitrag
\[
(-\lambda-m_{n-1})(-\lambda)^{n-1}
=
(-\lambda)^n - m_{n-1}(-\lambda)^{n-1}
=
-
(-1)^n
(\lambda^n + m_{n-1}\lambda^{n-1}).
\]
Zusammen mit \eqref{buch:komplex:fundamentalsatz:spur:eqn:unterdet}
folgt, dass die Determinante von $\bm{\alpha}-\lambda I$ das charakteristische
Polynom
\[
\chi_{\bm{\alpha}}(\lambda)
=
\det (\bm{\alpha}-\lambda I)
=
(-1)^n\biggl(
\lambda^n
+
\sum_{k=0}^{n-1} m_k \lambda^k
\biggr)
=
(-1)^n m(\lambda)
\]
ist.
Die Eigenwerte der Matrix sind also genau die Nullstellen des
ursprünglichen Minimalpolynoms $m(\lambda)$.

%
% Die Spur und die Determinante eines algebraischen Elementes
%
\subsubsection{Die Spur und die Determinante eines algebraischen Elements}
Die Multiplikation mit dem Element $\alpha$ wird in der Basis der
Potenzen von $\alpha$ durch die Matrix $\bm{\alpha}$ 
beschrieben.
Der Matrix $\bm{\alpha}$ kann man einige wohlbekannte Invarianten
zuordnen.

\begin{definition}[Spur und Determinante]
Die \emph{Spur} des algebraischen Elements $\alpha$ ist die Spur der
\index{Spur}%
\index{Determinante}%
Matrix $\bm{\alpha}$, die \emph{Determinante} von $\alpha$ ist die
Determinante der Matrix $\bm{\alpha}$:
\begin{equation}
\operatorname{tr}\alpha
=
\operatorname{tr}\bm{\alpha}
=
-m_{n-1}
\qquad\text{und}\qquad
\det\alpha
=
\det\bm{\alpha}
=
(-1)^nm_0.
\label{buch:komplex:symmetrisch:eqn:spurdet}
\end{equation}
\end{definition}

Die in der Definition enthaltene Aussage, dass Spur und Determinante
von $\alpha$ bereits durch die Koeffizienten des Minimalpolynoms bestimmt
sind, muss erst noch gerechtfertigt werden.
Die Spur und die Determinante einer Matrix sind Eigenschaften,
die nicht von der Wahl einer Basis abhängig sind.
Die Spur und die Determinante der linearen Abbildung
\[
\alpha\colon \mathbb{Q}(\alpha) \to \mathbb{Q}(\alpha) : x \mapsto \alpha x
\]
können also in irgend einer Basis berechnet werden.
Die Wahl von $\mathscr{B}=\{1,\alpha,\dots,\alpha^{n-1}\}$ als Basis
und die daraus resultierende spezielle Form der Matrix $\bm{\alpha}$,
führt auf die Werte \eqref{buch:komplex:symmetrisch:eqn:spurdet}.

\begin{korollar}
Sei $m(X)\in K[X]$ ein irreduzibles Polynom vom Grad $n$ mit Koeffizienten
in $K$.
Seien $\alpha_i$, $i=1,\dots,n$ die Nullstellen von $m(X)$.
Dann ist
\[
\operatorname{tr}\alpha_k
=
\sum_{i=1}^n \alpha_i
=
-m_{n-1}
\qquad\text{und}\qquad
\det \alpha_k
=
\prod_{i=1}^n \alpha_i
=
(-1)^nm_0
\]
für alle $k=1,\dots,n$.
\end{korollar}

\begin{proof}
Die Spur und die Determinante von $\alpha_k$ ist die Spur bzw.~die Determinante
der Matrix $\bm{\alpha_k}$.
Alle Nullstellen haben jedoch die gleiche, aus $m(X)$ konstruierte Matrix
\eqref{buch:komplex:fundamental:eqn:alphaMatrix},
sie haben also alle die gleiche Spur und Determinante, die ausserdem
die Summe bzw.~das Produkt der Eigenwerte dieser Matrix ist.
Das charakteristische Polynom der Matrix ist aber identisch mit dem
Minimalpolynom, die $\alpha_i$ sind somit auch die Eigenwerte.
\index{Eigenwert}%
\end{proof}

\begin{beispiel}
Sei $m(X) = X^3 + X + 1$.
Das Polynom $m(X)$ hat drei verschiedene Nullstellen.
Sei $\alpha$ eine dieser Nullstellen, dann ist
\[
\det \alpha = -1
\qquad\text{und}\qquad
\operatorname{tr}\alpha = 0.
\]
Tatsächlich kann man mit \texttt{octave} die numerischen Werte
\index{octave@\texttt{octave}}%
\[
\renewcommand{\arraycolsep}{2pt}
\begin{array}{rclcl}
\alpha_1 &=& \phantom{-}0.341163901914009 &+& 1.161541399997252i \\
\alpha_2 &=& \phantom{-}0.341163901914009 &-& 1.161541399997252i \\
\alpha_3 &=&           -0.682327803828019 & &
\end{array}
\]
finden.
Die Summe und das Produkt der Nullstellen sind
\[
\renewcommand{\arraycolsep}{2pt}
\begin{array}{rclcl}
\alpha_1+\alpha_2+\alpha_3 &=& -7.771561172376096\cdot 10^{-16} &\approx& -m_2,\\
\alpha_1\alpha_2\alpha_3   &=& -1                               &=&-m_0.
\end{array}
\]
Dies sind bis auf Maschinengenauigkeit die Koeffizienten des
Minimalpolynoms.
\end{beispiel}

%
% Symmetrische Funktionen der Nullstellen
%
\subsubsection{Symmetrische Funktionen der Nullstellen}
Die elementarsymmetrischen Polynome ausgewertet auf den Nullstellen
eines Polynoms mit Koeffizienten in $K$ ergeben Werte in $K$, selbst
dann wenn die Nullstellen selbst nicht in $K$ sind.
Zusammen mit dem
Fundamentalsatz~\ref{buch:komplex:fundamental:satz:fundamentalpolynom}
kann man daraus folgern, dass jedes symmetrische Polynom der Nullstellen
ebenfalls in $K$ liegt.

\begin{satz}
Sei $f(X_1,\dots,X_n)$ ein Polynom mit Koeffizienten in $K$,
das symmetrisch in den Variablen $X_1,\dots,X_n$ ist.
Sei weiter $p(X)\in K[X]$ ein Polynom mit den Nullstellen
$\alpha_1,\dots,\alpha_n$.
Dann ist $f(\alpha_1,\dots,\alpha_n)\in K$.
\end{satz}

\begin{proof}
Nach Satz~\ref{buch:komplex:fundamental:satz:fundamentalpolynom}
kann das Polynom $f(X_1,\dots,X_n)$ als Polynom $g(S_1,\dots,S_n)$
mit Koeffizienten in $K$ geschrieben werden, so dass
\[
f(X_1,\dots,X_n)
=
g(\sigma_1(X_1,\dots,X_n),\dots,\sigma_n(X_1,\dots,X_n))
\]
gilt.
Einsetzen der Nullstellen ergibt
\[
f(\alpha_1,\dots,\alpha_n)
=
g(
\underbrace{\sigma_1(\alpha_1,\dots,\alpha_n)}_{\displaystyle\in K}
,\dots,
\underbrace{\sigma_n(\alpha_1,\dots,\alpha_n)}_{\displaystyle\in K}
)
\in K.
\]
Es folgt, dass der Wert $f(\alpha_1,\dots,\alpha_n)$ in $K$ liegt.
\end{proof}

Wenn man zeigen kann, dass eine rationale Funktion von $X_1,\dots,X_n$
genau dann symmetrisch ist, wenn Zähler und Nenner symmetrisch sind,
dann folgt auch, dass die Werte rationaler, symmetrischer Funktionen
der Nullstellen im Koeffizientenkörper liegen.

\begin{satz}
Sei $f(X_1,\dots,X_n)$ eine rationale Funktion mit Koeffizienten
in $K$, die symmetrisch ist unter Vertauschungen der Argumente.
Dann kann $f$ geschrieben werden als Bruch
\[
f(X_1,\dots,X_n)
=
\frac{u(X_1,\dots,X_n)}{v(X_1,\dots,X_n)},
\]
wobei sowohl $u$ als auch $v$ symmetrische Polynome mit Koeffizienten
in $K$ sind.
Insbesondere ist $f(\alpha_1,\dots,\alpha_n)$ auch eine rationale Funktion
der Werte $\sigma_k(\alpha_1,\dots,\alpha_n)\in K$ mit Koeffizienten in $K$.
\end{satz}

\begin{proof}
Sei 
\[
f(X_1,\dots,X_n)
=
\frac{a(X_1,\dots,X_n)}{b(X_1,\dots,X_n)}
\]
mit Polynomen $a,b\in K[X_1,\dots,X_n]$.
Die rationale Funktion $f$ ist symmetrisch.
Insbesondere folgt
\begin{align*}
f(X_1,\dots,X_n)
&=
\frac{1}{n!}
\sum_{\pi\in S_n}
\frac{a^\pi(X_1,\dots,X_n)}{b^\pi(X_1,\dots,X_n)}.
\intertext{Durch Zusammenfassen entsteht ein Bruch der Form}
&=
\frac{
u(X_1,\dots,X_n)
}{
v(X_1,\dots,X_n)
},
\qquad
\text{mit }
v(X_1,\dots,X_n)
=
n!\prod_{\pi\in S_n} b^\pi(X_1,\dots,X_n).
\end{align*}
Der Nenner $v(X_1,\dots,X_n)$ ist offensichtlich ein symmetrisches
Polynom.
Da auch $f$ symmetrisch ist, folgt, dass das Polynom
\[
u(X_1,\dots,X_n)
=
f(X_1,\dots,X_n)\cdot v(X_1,\dots,X_n)
\]
ebenfalls symmetrisch ist.
Damit ist die Existenz der symmetrischen Polynome $u$ und $v$
gezeigt.
\end{proof}

\begin{korollar}
Sei $p(X)\in K[X]$ ein Polynom mit den Nullstellen
$\alpha_1,\dots,\alpha_n$ und
$f(X_1,\dots,X_n)$ eine rationale Funktion mit Koeffizienten
in $K$, die symmetrisch ist unter Vertauschungen der Argumente.
Dann ist $f(\alpha_1,\dots,\alpha_n)$ auch eine rationale Funktion
der Werte $\sigma_k(\alpha_1,\dots,\alpha_n)\in K$ mit Koeffizienten in $K$.
\end{korollar}

Das Korollar bedeutet, dass man die Nullstellen eines Polynoms 
nicht kennen muss, wenn man die Werte einer symmetrischen rationalen
Funktion der Nullstellen berechnen will, weil diese Werte bereits vollständig
durch die Koeffizienten des Polynoms bestimmt sind.
